\begin{abstract}
コヒーレント・イジングマシン（CIM）の出力を後段の探索の初期解（種）に使うとき，同じカット値の種は同じ価値を持つか，その差を安価に予測して選別すると役に立つかを，MAX-CUT（G-set と全列挙できる小規模問題）で調べた．種の価値は，その種から始めた焼きなまし法（SA）の最良値の期待値とした．(1) カット値と 1 反転局所探索の改善量で説明できない価値の差は，SA が種の近くを探す設定では，CIM の出力が多様な 7 問題すべてで再現した（split-half 信頼性 0.67〜0.98）．SA が種から大きく離れる設定では差は消えた．(2) 探索的に見つけた固有ベクトル指標（CIM の線形増幅で最初に成長するモードへの揃い方）を事前登録して検証し，G22 で再現した（$\rho=0.56$，95\% 区間 [0.46, 0.65]）．G1 でも同じ大きさで成り立ったが，ほかの問題では効かなかった．全列挙した小規模問題では，価値は最適解までの距離でほぼ決まり（層内 Kendall $\tau=-0.74$），エネルギー障壁の効果は弱かった．問題をまたいで同じ向きに効いた指標は，短い SA の試し打ちだけだった．(3) CIM の生成を含めた同じ総時間で比べると，カット値を超える選別は最良値をほとんど改善しなかった（最大 0.003\%）．CIM の種は，後段の予算が短い（500 掃引）とランダムな初期解より良かったが，長い（5000 掃引）と CIM の時間を SA に回したほうが良かった．(4) 後段がメメティックアルゴリズム（GA）の場合は，CIM の生成時間を含めても，CIM の初期集団がランダムな初期集団より良かった（G55・G70）．ただし，ハミング距離で多様な初期集団を選ぶと，カット値だけで選ぶより悪くなった．付随して，GA のタブー探索の同点処理の不具合を見つけて修正した．以前の「CIM の初期集団は GA に役立たない」という結論は，この不具合によるものだった．
\end{abstract}
